% Audit de veracite des citations de rapport de jury - 2026-09-20
% Bloc 4 : citation attribuee au jury 2022 introuvable dans le corpus ; le theme
% (cablage API / variateur) figure en 2021. Remplacee par le verbatim du
% Rapport_Jury_2021 (l.819), annee corrigee en 2021.
\expandafter\ifx\csname CAPETINCLUDEMODE\endcsname\relax
\documentclass[11pt,a4paper]{article}
% preamble.tex — packages et commandes communes a toutes les fiches CAPET SII IE
\input{glyphtounicode}
\pdfgentounicode=1
\usepackage{cmap}
\usepackage[T1]{fontenc}
\usepackage[utf8]{inputenc}
\usepackage[french]{babel}
% babel[french] active a \begin{document} des raccourcis de ponctuation (:;!?)
% qui cassent la syntaxe to[...] de circuitikz (erreur "+ or - expected") ;
% on les desactive, ce qui ne change que l'espacement fin avant ces signes.
\AtBeginDocument{\shorthandoff{:;!?}}
\usepackage[a4paper, top=1.6cm, bottom=1.6cm, left=1.4cm, right=1.4cm]{geometry}
\setlength{\headheight}{14.5pt}
\usepackage{amsmath, amssymb, amsfonts}
\usepackage{siunitx}
\sisetup{locale=FR, output-decimal-marker={,}, inter-unit-product=\ensuremath{{}\cdot{}},
         group-minimum-digits=4}
% Unites non fournies par siunitx, declarees ici une fois pour toutes (fiches + exercices).
% Les gardes \ifdefined evitent tout conflit avec une redeclaration locale restee dans un
% fichier : ne pas retirer les gardes, ne pas redeclarer ces unites dans les fichiers.
\DeclareSIUnit{\jour}{j}
\ifdefined\var\else\DeclareSIUnit{\var}{var}\fi   % puissance reactive
\ifdefined\vad\else\DeclareSIUnit{\vad}{VAD}\fi   % puissance deformante
\ifdefined\tr\else\DeclareSIUnit{\tr}{tr}\fi      % tour (tr/min)
\usepackage{array, booktabs, multirow, tabularx}
\usepackage{xcolor}
\usepackage[most, breakable]{tcolorbox}
\usepackage{multicol}
\usepackage{enumitem}
\usepackage{mathtools}
\usepackage{lmodern}
% Refonte visuelle 2026 (direction "moderne coloré") : Lato pour le texte,
% maths sans-serif assorties via newtxsf (charge apres amsmath/amssymb).
\usepackage[default]{lato}
\usepackage{newtxsf}
\usepackage{microtype}
\usepackage{fancyhdr}
\usepackage{tikz}
\usetikzlibrary{positioning,arrows.meta}
\usepackage{circuitikz}
\usepackage{lastpage}

% ── Couleurs par domaine (palette modernisée 2026) ────────────────────
\definecolor{colElec}{RGB}{29, 78, 216}        % Électrotechnique — bleu vif
\definecolor{colElecLight}{RGB}{219, 234, 254} % fond bleu clair
\definecolor{colAuto}{RGB}{21, 128, 61}        % Automatique — vert
\definecolor{colMeca}{RGB}{194, 65, 12}        % Mécanique — orange
\definecolor{colTherm}{RGB}{185, 28, 28}       % Thermique — rouge
\definecolor{colQT}{RGB}{232, 244, 233}        % (hérité, plus utilisé par questiontype)
\definecolor{colJury}{RGB}{255, 248, 225}      % fond jury (crème)
\definecolor{colJuryFrame}{RGB}{202, 138, 4}   % bordure/bandeau jury (ambre)
\definecolor{colGrayLight}{RGB}{245, 245, 247} % fond gris neutre
\definecolor{colSeparator}{RGB}{180, 180, 180} % gris séparateur
\definecolor{colExo}{RGB}{8, 145, 178}          % Banque d'exercices — cyan (distinct des fiches)
% ── Encre et gris de la refonte ───────────────────────────────────────
\definecolor{colInk}{RGB}{23, 32, 42}          % texte courant (quasi-noir)
\definecolor{colSlate}{RGB}{83, 98, 118}       % gris ardoise (méta, en-têtes)
\definecolor{colRule}{RGB}{203, 213, 225}      % filets clairs
\definecolor{colJuryBg}{RGB}{255, 248, 225}    % fond crème jury
\definecolor{colPresence}{RGB}{83, 98, 118}    % bloc 5 (présence dans les sujets)
\AtBeginDocument{\color{colInk}}

% ── Style encadré exercice + corrigé (banque d'exercices) ─────────────
\tcbset{
  exoBox/.style={
    breakable,
    enhanced,
    colback=colExo!8!white,
    colframe=colExo,
    colbacktitle=colExo,
    coltitle=white,
    fonttitle=\bfseries\small,
    left=4pt, right=4pt, top=3pt, bottom=3pt,
    before skip=5pt, after skip=5pt,
  }
}

% ── Bascule enonce / corrige (banque d'exercices) ─────────────────────
% \CAPETCORRIGE est defini en ligne de commande (pdflatex "\def\CAPETCORRIGE{1}\input{...}")
% pour compiler la variante corriges d'un meme fichier source ; par defaut (non defini),
% la compilation produit la variante enonces. \CAPETSHOW{enonce}{corrige} choisit la branche.
%
% TROIS MODES (refonte 2026-09-20, manuel de revision) :
%   - aucun drapeau   -> imprime l'enonce            (variante enonces, INCHANGE)
%   - \CAPETCORRIGE   -> imprime le corrige          (variante corriges, INCHANGE)
%   - \CAPETMANUEL    -> imprime l'enonce ET met le corrige de cote pour l'annexe
%
% Le mode manuel stocke chaque corrige dans une macro (\csgdef), sans ecriture
% fichier ni re-tokenisation : le contenu (circuitikz, accents, %) traverse intact.
% L'annexe les rejoue ensuite dans l'ordre via \CAPETANNEXE.
\usepackage{etoolbox}

\newcounter{CAPETexo}          % numero continu d'exercice dans le manuel
\setcounter{CAPETexo}{0}
\newcounter{CAPETannexe}       % compteur de relecture en annexe

% Titre courant de la notion en cours (pose par \CAPETNOTION), sert d'intertitre annexe
\newcommand{\CAPETnotioncode}{}
\newcommand{\CAPETnotiontitre}{}

% \CAPETNOTION{code}{titre} : ouvre une notion dans le manuel (memorise le libelle,
% et marque le rang du premier exercice de la notion pour l'intertitre d'annexe).
\newcommand{\CAPETNOTION}[2]{%
  \renewcommand{\CAPETnotioncode}{#1}%
  \renewcommand{\CAPETnotiontitre}{#2}%
  \csgdef{CAPETnotionstart\the\numexpr\value{CAPETexo}+1\relax}{#1 -- #2}%
}

% ── \degree robuste ───────────────────────────────────────────────────
% Certains fichiers d'exercices utilisent \degree. Ni le preambule ni siunitx
% ne le definissent, et \providecommand local ne suffit pas : sous hyperref,
% \degree existe deja comme commande factice (\ERROR), donc \providecommand
% s'abstient et l'appel echoue a la compilation du manuel (erreur invisible
% en compilation isolee). On le (re)definit donc inconditionnellement ici,
% apres le chargement des packages, pour tout le corpus.
\AtBeginDocument{\providecommand{\degree}{}\renewcommand{\degree}{\ensuremath{^\circ}}}

% Libelle decoratif : les fichiers d'exercices utilisent aussi \CAPETSHOW dans
% leur \fancyhead et leur bloc de titre, pour afficher "(enonces)" ou "(corriges)".
% Ces appels-la ne sont PAS des exercices : en mode manuel ils ne doivent ni
% consommer un numero, ni alimenter l'annexe. On les reconnait a leur premier
% argument, strictement egal au libelle ci-dessous.
\newcommand{\CAPETLIBELLEENONCES}{(\'enonc\'es)}

\ifdefined\CAPETMANUEL
  % --- Mode manuel : enonce en place, corrige differe vers l'annexe ---
  \newcommand{\CAPETSHOW}[2]{%
    \def\CAPETargun{#1}%
    \ifx\CAPETargun\CAPETLIBELLEENONCES
      #1%            appel decoratif : rendu tel quel, hors numerotation
    \else
      \CAPETSHOWEXO{#1}{#2}%
    \fi
  }
  \newcommand{\CAPETSHOWEXO}[2]{%
    \stepcounter{CAPETexo}%
    \csgdef{CAPETcorrige\theCAPETexo}{#2}%
    \edef\CAPETtmpnum{\theCAPETexo}%
    #1%
    \par\vspace{2pt}%
    {\footnotesize\itshape\color{colExo}%
     \hyperref[capet:cor:\CAPETtmpnum]{Corrige de l'exercice \CAPETtmpnum\ $\rightarrow$ p.~\pageref{capet:cor:\CAPETtmpnum}}%
     \label{capet:eno:\CAPETtmpnum}}%
  }
  % \CAPETANNEXE : rejoue tous les corriges captures, dans l'ordre, avec intertitres
  \newcommand{\CAPETANNEXE}{%
    \setcounter{CAPETannexe}{0}%
    \loop\ifnum\value{CAPETannexe}<\value{CAPETexo}%
      \stepcounter{CAPETannexe}%
      \edef\CAPETcurnum{\theCAPETannexe}%
      \ifcsname CAPETnotionstart\CAPETcurnum\endcsname
        \par\medskip
        {\normalsize\bfseries\color{colExo}\csuse{CAPETnotionstart\CAPETcurnum}}%
        \par\nobreak\smallskip
      \fi
      \begin{tcolorbox}[exoBox, breakable,
        title={Corrige de l'exercice \CAPETcurnum\
               \hfill\normalfont\footnotesize
               \hyperref[capet:eno:\CAPETcurnum]{enonce p.~\pageref{capet:eno:\CAPETcurnum}}}]%
      \label{capet:cor:\CAPETcurnum}%
      \csuse{CAPETcorrige\CAPETcurnum}%
      \end{tcolorbox}%
    \repeat
  }
\else
  % --- Modes historiques, comportement strictement inchange ---
  \newcommand{\CAPETSHOW}[2]{\ifdefined\CAPETCORRIGE #2\else #1\fi}
  \newcommand{\CAPETANNEXE}{}
\fi

% ── Difficulté 1 à 3 (étoiles pleines, amssymb) ───────────────────────
\newcommand{\difficulte}[1]{%
  \ifcase#1\relax
  \or $\bigstar$\;niveau 1/3%
  \or $\bigstar\bigstar$\;niveau 2/3%
  \or $\bigstar\bigstar\bigstar$\;niveau 3/3%
  \fi
}

% ── Style encadré extrait de document technique (banque d'exercices) ──
\tcbset{
  docTech/.style={
    breakable,
    enhanced,
    colback=colGrayLight,
    colframe=colSeparator,
    fonttitle=\bfseries\small\itshape,
    coltitle=black,
    colbacktitle=colGrayLight,
    left=4pt, right=4pt, top=3pt, bottom=3pt,
    before skip=4pt, after skip=4pt,
  }
}

% ── Style tcolorbox de section principale ────────────────────────────
% Refonte 2026 : bandeau titre coloré arrondi + fond très pâle, sans cadre.
\tcbset{
  ficheSection/.style={
    breakable, enhanced jigsaw,
    frame hidden, boxrule=0pt, arc=1.8mm,
    colback=#1!4!white,
    colbacktitle=#1,
    coltitle=white,
    fonttitle=\sffamily\bfseries,
    left=7pt, right=7pt, top=5pt, bottom=5pt,
    toptitle=3pt, bottomtitle=3pt,
    before skip=10pt, after skip=8pt,
  }
}

% ── Style encadré question type ──────────────────────────────────────
% Carte blanche, bandeau sombre, liseré très clair.
\tcbset{
  questiontype/.style={
    breakable, enhanced jigsaw,
    colback=white,
    frame hidden, boxrule=0pt, arc=1.8mm,
    colbacktitle=colElec!78!black,
    coltitle=white,
    fonttitle=\sffamily\bfseries\small,
    left=7pt, right=7pt, top=4pt, bottom=4pt,
    toptitle=2.5pt, bottomtitle=2.5pt,
    before skip=8pt, after skip=8pt,
    borderline={0.6pt}{0pt}{colElec!12!white},
  }
}

% ── Style encadré jury ───────────────────────────────────────────────
% Bandeau ambre, fond crème, sans cadre.
\tcbset{
  jury/.style={
    breakable, enhanced jigsaw,
    colback=colJuryBg,
    frame hidden, boxrule=0pt, arc=1.8mm,
    colbacktitle=colJuryFrame,
    coltitle=white,
    fonttitle=\sffamily\bfseries,
    left=7pt, right=7pt, top=4pt, bottom=4pt,
    toptitle=3pt, bottomtitle=3pt,
    before skip=10pt, after skip=8pt,
  }
}

% ── En-tête / pied de page communs (refonte 2026) ─────────────────────
% #1 code fiche, #2 titre court, #3 couleur domaine
\newcommand{\ficheHeaders}[3]{%
  \fancyhead[L]{\small\sffamily{\bfseries\textcolor{#3}{#1}}\hspace{0.7em}\textcolor{colSlate}{#2}}%
  \fancyhead[R]{\small\sffamily\textcolor{colSlate}{CAPET SII IE\, --- \,Session 2026}}%
  \fancyfoot[C]{\small\sffamily\textcolor{colSlate}{\thepage\,/\,\pageref{LastPage}}}%
  \fancyhead[C]{}\fancyfoot[L]{}\fancyfoot[R]{}%
}

% ── Bloc en-tête de fiche (BLOC 0, refonte 2026) ──────────────────────
% #1 code  #2 titre  #3 couleur domaine  #4 domaine  #5 ligne présence  #6 sous-thèmes
\newcommand{\ficheEntete}[6]{%
  \begin{tcolorbox}[enhanced, breakable,
    frame hidden, boxrule=0pt, arc=2.5mm,
    interior style={left color=#3, right color=#3!75!black},
    left=9pt, right=9pt, top=7pt, bottom=7pt,
    before skip=0pt, after skip=6pt]
    {\color{white}\sffamily\bfseries\LARGE #1\quad #2}\\[4pt]
    {\color{#3!22!white}\sffamily\small \MakeUppercase{#4}\quad\textbar\quad #5}\\[1pt]
    {\color{#3!35!white}\sffamily\footnotesize Sous-thèmes : #6}
  \end{tcolorbox}%
}

% ── Mise en page ─────────────────────────────────────────────────────
\pagestyle{fancy}
\fancyhf{}
\renewcommand{\headrulewidth}{0.5pt}
\renewcommand{\footrulewidth}{0pt}
\AtBeginDocument{%
  \renewcommand{\headrule}{{\color{colRule}\hrule width\headwidth height 0.5pt}}%
}

\setlist[itemize]{leftmargin=1.2em, itemsep=1pt, topsep=2pt, parsep=0pt}
\setlist[enumerate]{leftmargin=1.2em, itemsep=3pt, topsep=2pt, parsep=0pt}

\setlength{\parskip}{2pt}
\setlength{\columnsep}{10pt}

\def\CAPETBEGINFILE{\begin{document}}
\def\CAPETENDFILE{\end{document}}
\else
\def\CAPETBEGINFILE{}
\def\CAPETENDFILE{}
\fi

\ficheHeaders{T29}{Logique séquentielle, API}{colAuto}

\CAPETBEGINFILE

% -----------------------------------------------------------------------
%  BLOC EN-TETE
% -----------------------------------------------------------------------

\ficheEntete{T29}{Logique séquentielle, GRAFCET, automates programmables}{colAuto}%
  {Automatique / électronique numérique - Séquentiel}%
  {\textbf{Présence :} \textbf{5/13} sujets (2013-2026) - thème \emph{occasionnel}}%
  {bascules RS/D/JK - GRAFCET (au programme, jamais interrogé 2013-2026) - langages IEC 61131-3 - câblage API + variateur - algorigramme - protocoles série (UART, liaison GPS)}

% -----------------------------------------------------------------------
%  1. RAPPELS DE COURS
% -----------------------------------------------------------------------

\begin{tcolorbox}[
  ficheSection=colAuto,
  title={1. Rappels de cours},
]

% -- 1.1 Bascules -----------------------------------------------------
\textbf{1.1 Bascules (éléments mémorisants)}

\begin{center}
\renewcommand{\arraystretch}{1.2}
\begin{tabular}{lccl}
\toprule
\textbf{Bascule} & \textbf{Entrée(s)} & \textbf{$Q_{n+1}$} & \textbf{Remarque} \\
\midrule
RS & $S=0,R=0$ & $Q_n$ & Mémorisation \\
   & $S=1,R=0$ & 1 & Set \\
   & $S=0,R=1$ & 0 & Reset \\
   & $S=1,R=1$ & Interdit & - \\
\midrule
D  & $D$ & $D$ & Sur front d'horloge \\
JK & $J=0,K=0$ & $Q_n$ & Mémorisation \\
   & $J=1,K=1$ & $\overline{Q_n}$ & Toggle \\
T  & $T$ & $Q_n \oplus T$ & Toggle si $T=1$ \\
\bottomrule
\end{tabular}
\end{center}

\medskip\hrule\vspace{4pt}

% -- 1.2 GRAFCET ------------------------------------------------------
\textbf{1.2 GRAFCET (IEC 60848)}

\emph{N.B. — Dans les 41 sujets et rapports analysés (2013-2026), le mot « GRAFCET » n'apparaît que deux fois (référence au programme dans le sujet commun 2017), jamais comme question d'examen : les questions d'automatisme portent sur l'algorigramme, le Texte Structuré et le câblage API/variateur. Le GRAFCET reste néanmoins au programme comme langage de description (graphe d'états, logigramme, algorigramme).}

\smallskip
\textbf{Éléments constitutifs :}
\begin{itemize}
  \item \textbf{Étape} : rectangle numéroté (double contour = étape initiale) ; actions associées
  \item \textbf{Transition} : trait horizontal avec réceptivité (condition booléenne)
  \item \textbf{Liaison orientée} : flèche (sens de haut en bas par convention)
\end{itemize}

\smallskip
\textbf{Extrait de GRAFCET (séquence linéaire) :}
\begin{center}
\begin{tikzpicture}[>=Stealth, font=\small]
  \node[draw, double, minimum width=8mm, minimum height=8mm] (e0) at (0,4) {0};
  \draw (0,3.6) -- (0,3.2);
  \draw[thick] (-0.3,3.2) -- (0.3,3.2);
  \node[right] at (0.4,3.2) {$dcy$};
  \draw[->] (0,3.2) -- (0,2.6);
  \node[draw, minimum width=8mm, minimum height=8mm] (e1) at (0,2.2) {1};
  \draw (0,1.8) -- (0,1.4);
  \draw[thick] (-0.3,1.4) -- (0.3,1.4);
  \node[right] at (0.4,1.4) {$t_1$};
  \draw[->] (0,1.4) -- (0,0.8);
  \node[draw, minimum width=8mm, minimum height=8mm] (e2) at (0,0.4) {2};
\end{tikzpicture}
\end{center}
L'étape 0 (double contour) est active au repos. La transition franchie dès que $dcy$ (dcy = mise en marche) est vraie active l'étape 1 ; puis $t_1$ franchie active l'étape 2.

\textbf{Règles d'évolution :}
\begin{enumerate}
  \item L'étape initiale est active au démarrage
  \item Une transition est franchissable si \textbf{toutes} les étapes amont sont actives ET sa réceptivité est vraie
  \item Le franchissement désactive les étapes amont et active les étapes aval
  \item Plusieurs transitions franchissables le même cycle le sont toutes simultanément
\end{enumerate}

\textbf{Structures courantes :}

\begin{center}
\renewcommand{\arraystretch}{1.15}
\begin{tabular}{ll}
\toprule
\textbf{Structure} & \textbf{Représentation} \\
\midrule
Séquence simple & Étapes en série \\
Divergence OU & Plusieurs branches alternatives (aiguillage) \\
Convergence OU & Réunion de branches alternatives \\
Divergence ET & Double barre - branches simultanées \\
Convergence ET & Double barre de synchronisation \\
Saut / reprise & Liaison directe entre étapes non consécutives \\
\bottomrule
\end{tabular}
\end{center}

\medskip\hrule\vspace{4pt}

% -- 1.3 Langages IEC 61131-3 -----------------------------------------
\textbf{1.3 Langages de programmation automate (IEC 61131-3)}

\begin{center}
\renewcommand{\arraystretch}{1.15}
\begin{tabular}{lll}
\toprule
\textbf{Langage} & \textbf{Type} & \textbf{Usage typique} \\
\midrule
Ladder (LD) & Graphique & Logique contacts/bobines \\
FBD & Graphique & Blocs PID, temporisateurs \\
Texte Structuré (ST) & Textuel & Algorithmie, calculs \\
SFC (GRAFCET) & Graphique & Séquencement \\
\bottomrule
\end{tabular}
\end{center}

\textbf{Temporisateurs :} TON (ON-delay), TOF (OFF-delay), TP (impulsion). TON : sortie $Q$ active uniquement si l'entrée reste vraie pendant toute la durée $PT$.

\medskip\hrule\vspace{4pt}

% -- 1.4 Ceblage API --------------------------------------------------
\textbf{1.4 Câblage automate + variateur de vitesse}

\begin{itemize}
  \item \textbf{Sortie analogique API} (0-10 V ou 4-20 mA) - borne de consigne variateur (AI1)
  \item \textbf{Sorties TOR API} - LI1 (sens avant), LI2 (sens arrière), LI3 (2\textsuperscript{e} vitesse)
  \item \textbf{Entrées TOR API} - SD (signal défaut), RUN, disponible
  \item \textbf{Entrée analogique API} - retour tachymétrique ou sortie AO variateur
  \item \textbf{Commun 24 V DC} : toujours câbler le 0 V de référence des entrées logiques
\end{itemize}

\medskip\hrule\vspace{4pt}

% -- 1.5 Protocoles serie ----------------------------------------------
\textbf{1.5 Liaison série UART}

Format typique : 1 bit start + 8 bits data + 1 bit stop = \textbf{10 bits par caractère}.

\[
  T_{trame} = \frac{N_{car} \times 10}{f_{bauds}}
  \qquad
  \textit{Ex. (2023) : } \frac{71 \times 10}{115\,200} \approx \SI{6{,}16}{\milli\second}
\]

\end{tcolorbox}

% -----------------------------------------------------------------------
%  2. FORMULAIRE SYNTHETIQUE
% -----------------------------------------------------------------------

\begin{tcolorbox}[
  ficheSection=colAuto,
  title={2. Formulaire synthétique},
]
\small
\begin{multicols}{2}
\renewcommand{\arraystretch}{1.45}
\begin{tabular}{@{}p{5.2cm}l@{}}
\toprule
\textbf{Formule / Règle} & \textbf{Usage} \\
\midrule
$Q_{n+1}^D = D$ & Bascule D \\
$Q_{n+1}^T = Q_n \oplus T$ & Bascule T (toggle) \\
Transition franchissable & Étapes amont actives ET réceptivité vraie \\
TON : $Q$ actif si IN vrai pendant $PT$ & Temporisateur \\
\bottomrule
\end{tabular}

\columnbreak

\renewcommand{\arraystretch}{1.45}
\begin{tabular}{@{}p{5.2cm}l@{}}
\toprule
\textbf{Formule / Règle} & \textbf{Usage} \\
\midrule
$I_{th} \approx 0{,}94 \times I_n$ (moteur) & Choix variateur \\
Plage 4-20 mA : 0 mA = défaut câblé & Sécurité \\
$T_{trame} = N_{car} \times 10 / f_{bauds}$ & UART \\
Affectation ST : $:=$, comparaison $=$ & Syntaxe \\
\bottomrule
\end{tabular}
\end{multicols}
\normalsize
\end{tcolorbox}

% -----------------------------------------------------------------------
%  3. QUESTIONS TYPES
% -----------------------------------------------------------------------

\begin{tcolorbox}[
  ficheSection=colAuto,
  title={3. Questions types - extraites des sujets réels},
]

% -- QT1 -------------------------------------------------------------

\begin{tcolorbox}[questiontype,
  title={QT1 - Câblage API + variateur ATV630 \hfill\normalfont\footnotesize E2 appl. commune 2021 (Q.19-20)}
]
\textbf{Énoncé.}
Compléter les schémas de puissance et de commande du variateur Schneider ATV630 (\SI{90}{kW}). Câbler le débitmètre, le capteur de pression, l'électrovanne et le pressostat aux modules d'E/S du Schneider M340.

\tcbline

\textbf{Corrigé — éléments clés attendus.}
\begin{itemize}
  \item Schéma puissance : alimentation triphasée - disjoncteur-moteur - bornes L1/L2/L3 du variateur ; bornes U/V/W - moteur. Courant variateur : $I_{th} \geq 0{,}94 \times I_{nom} \approx \SI{170}{\ampere}$.
  \item Schéma commande : sortie AO du M340 (0-10 V ou 4-20 mA) - AI1 du variateur (consigne vitesse).
  \item Sortie TOR M340 - LI1 (marche/arrêt) ; LI2 si inversion de sens.
  \item Entrée TOR M340 - SD (défaut variateur) pour signalisation.
  \item \textbf{Ne pas oublier :} commun 0 V entre l'API et le variateur (borne COM).
\end{itemize}

\textbf{Conclusion.} L'architecture de puissance et de commande, ainsi que les points de câblage identifiés (consigne, retour défaut, commun 0 V), sont conformes au cahier des charges fonctionnel de la liaison API-variateur.
\end{tcolorbox}

\vspace{4pt}

% -- QT2 -------------------------------------------------------------

\begin{tcolorbox}[questiontype,
  title={QT2 - Algorigramme de gestion de volume \hfill\normalfont\footnotesize E2 appl. commune 2021 (Q.24)}
]
\textbf{Énoncé.}
Proposer un algorigramme du bloc fonction « Contrôler le volume ». Définir les sorties \textit{Res\_plein} et \textit{Res\_vide} en fonction de la mesure de volume.

\tcbline

\textbf{Corrigé.}

\textbf{Étape 1 — Entrée.} $V_{mes}$ [m$^3$] (mesure de volume issue de la chaîne capteur + CAN).

\textbf{Étape 2 — Tests de seuil.}
\begin{enumerate}
  \item SI $V_{mes} \geq 0{,}95 \times V_{max}$ - \textit{Res\_plein} $:= 1$
  \item SI $V_{mes} \leq 0{,}05 \times V_{max}$ - \textit{Res\_vide} $:= 1$
  \item Sinon : \textit{Res\_plein} $:= 0$ ; \textit{Res\_vide} $:= 0$
\end{enumerate}

Structure de contrôle : deux SI-ALORS-SINON imbriqués (ou juxtaposés). En Texte Structuré :
\begin{center}
\small\ttfamily
IF Vmes >= 0.95*Vmax THEN Res\_plein := TRUE; ELSE Res\_plein := FALSE; END\_IF;
\end{center}

\textbf{Conclusion.} Cette logique de seuils à 5\,\% de marge de part et d'autre évite les basculements intempestifs autour des bornes exactes : elle est directement exploitable en API et conforme au cahier des charges de surveillance du volume.
\end{tcolorbox}

\vspace{4pt}

% -- QT3 -------------------------------------------------------------

\begin{tcolorbox}[questiontype,
  title={QT3 - Durée de transmission UART \hfill\normalfont\footnotesize E2 disciplinaire 2023}
]
\textbf{Énoncé.}
Une trame GPS (protocole NMEA) comporte 71 caractères. La liaison est configurée à 115 200 bauds, format 8N1. Calculer la durée de transmission de la trame.

\tcbline

\textbf{Corrigé.}

\textbf{Étape 1 — Bits par caractère.} Format 8N1 : 1 bit start + 8 bits data + 1 bit stop = 10 bits/caractère.

\textbf{Étape 2 — Durée de trame.}
\[
  T_{trame} = \frac{71 \times 10}{115\,200} \approx \SI{6{,}16}{\milli\second}
\]

\textbf{Conclusion.} La tâche automate qui lit la liaison série doit avoir une période de scrutation strictement inférieure à $T_{trame} \approx 6{,}16\,\mathrm{ms}$ pour éviter la perte de trame : c'est cette contrainte de scrutation, et non le seul calcul de durée, qui conditionne la conformité de l'architecture de communication.
\end{tcolorbox}

\end{tcolorbox}

% -----------------------------------------------------------------------
%  4. ATTENTES DU JURY
% -----------------------------------------------------------------------

\begin{tcolorbox}[
  jury,
  title={4. Attentes du jury - Erreurs récurrentes et points de vigilance},
]

\begin{enumerate}[label=\textbf{\arabic*.}, leftmargin=*, itemsep=5pt]

  \item \textbf{Câblage variateur - le commun 24 V.} Oublier le câblage du commun 0 V entre l'API et le variateur est la première cause d'erreur. Les entrées logiques LI du variateur sont référencées à ce commun.

  \item \textbf{4-20 mA : 0 mA signifie une rupture de câble.} C'est l'avantage de la boucle de courant sur le 0-10 V : 0 mA est détectable comme un défaut.

  \item \textbf{Algorigramme : structures de contrôle.} Le jury déplore que les "notions de base d'algorithmie ne soient pas maîtrisées" (2021). Distinguer SI (test unique), TANT QUE (boucle conditionnelle), POUR (boucle à compteur).

  \item \textbf{Langage Texte Structuré.} L'affectation est $:=$ (pas $=$) ; la comparaison est $=$. Terminer chaque instruction par un point-virgule.

  \item \textbf{GRAFCET - action vs réceptivité.} L'action est associée à l'étape (ce qui se fait pendant que l'étape est active) ; la réceptivité est associée à la transition (condition pour quitter l'étape).

  \item \textbf{Schéma électrique : point noir récurrent.} Moyenne de 4,8/20 en 2021 sur la partie câblage. Le jury attend la maîtrise complète du schéma variateur (puissance + commande).

\end{enumerate}

\medskip
\begin{tcolorbox}[
  colback=colJuryFrame!12!colJury,
  colframe=colJuryFrame,
  left=4pt, right=4pt, top=2pt, bottom=2pt,
  fontupper=\small\itshape
]
Rapport jury 2021 : «~Les candidats doivent acquérir des connaissances dans l'algorithmie et dans la réalisation de schémas électriques prenant appui sur de la documentation constructeur. Ces notions sont largement accessibles à un étudiant de BTS.~»\\
Rapport jury 2021 : «~Le jury constate que beaucoup de candidats ne sont pas à l'aise dans la réalisation de schémas électriques. Un candidat au CAPET doit être capable de câbler une entrée ou une sortie TOR analogique d'un automate industriel programmable, ainsi que la partie commande d'un variateur de fréquence.~»
\end{tcolorbox}

\end{tcolorbox}

% -----------------------------------------------------------------------
%  PIED - PRESENCE DANS LES SUJETS
% -----------------------------------------------------------------------

\begin{tcolorbox}[
  ficheSection=colPresence,
  title={Présence dans les sujets analysés (2013-2026)},
]
\begin{center}
\footnotesize\renewcommand{\arraystretch}{1.15}
\begin{tabular}{lccccccccccccc}
\toprule
\textbf{Année} & 13 & 14 & 15 & 16 & 17 & 18 & 19 & 20 & 21 & 22 & 23 & 24 & 25 \\
\midrule
Câblage API/variateur & - & - & - & - & - & $\checkmark$ & - & - & $\checkmark$ & $\checkmark$ & - & - & - \\
Algorigramme / ST   & - & - & - & - & - & - & - & - & $\checkmark$ & - & - & $\checkmark$ & - \\
Protocoles série    & - & - & - & - & - & - & - & - & - & - & $\checkmark$ & $\checkmark$ & - \\
\bottomrule
\end{tabular}
\end{center}
\smallskip
\footnotesize\textbf{Sous-thèmes fréquents :}
câblage API + variateur (puissance + commande) - algorigramme bloc fonction - choix du variateur ($I_{th}$) - protocole UART (durée trame, chronogramme) - algorithmie en Texte Structuré

\medskip
\textbf{Ordres de grandeur (repères rapides).}
\begin{itemize}[leftmargin=1.3em]
  \item Réseau BT FR : tension phase-neutre typique \SI{230}{\volt}, tension phase-phase typique \SI{400}{\volt}, fréquence \SI{50}{\hertz}.
  \item Courants d'étude fréquents en sujet : de quelques \si{\ampere} à quelques dizaines de \si{\ampere}.
  \item En conversion d'énergie, un rendement attendu reste généralement entre 0{,}8 et 0{,}98 selon la technologie.
\end{itemize}

\medskip
\textbf{Synthèse express (à refaire sans support).}
\begin{enumerate}[leftmargin=1.4em]
  \item Poser les hypothèses du modèle et les conventions de signe.
  \item Écrire 2 à 4 relations clés avec unités.
  \item Réaliser un calcul numérique complet et vérifier l'ordre de grandeur.
  \item Justifier le résultat par un argument technique de niveau CAPET.
\end{enumerate}

\end{tcolorbox}

\CAPETENDFILE
